\section{Discussion}
\label{sec:discussion}
The evidence favours matching generation to the intended task. Scene-based
methods offer scenario and label control; learned methods fit scan
regularities or complete observations; hybrids combine these strengths
while retaining real-data requirements. PreSIL, LiDARsim and SynLiDAR
demonstrate transfer in different protocols, whereas LiDARGen's completion
results cannot establish synthetic-training gains. These distinctions
prevent a misleading cross-paper ranking.

Realism and useful variation are also different objectives. Improving sensor
sampling can help within a pipeline, but LiDARGen's metric disagreements and
the paired-scenario results undermine a universal realism-to-utility rule.
Rare-case coverage deserves separate scrutiny: matching common observations
may omit the very cases synthesis is intended to add.
\Textcite{shumailovAIModelsCollapse2024} show loss of distribution tails in
recursive training on generated data. Their language, VAE and mixture-model
settings motivate preserving real-data provenance, but do not establish
collapse in one-pass LiDAR augmentation.

A focused next study could vary ray drop or angular sparsity while fixing
scenes, labels, model and training budget, then compare diagnostic changes
with class- and range-specific real performance. This connects existing
sensor and transfer evaluations rather than claiming evaluation methods
are absent; its novelty still requires broader searching.

The review's selected collection, incomplete systematic screening and
pending author verification limit generality. Sensor, label and protocol
differences preclude pooled effects. Compute costs and uncertainty have
not been systematically extracted, and reported validation gains are not
guarantees of deployment performance.
